% Preámbulo común de la memoria y de la guía de uso.
% Tipografía: Libertinus (texto y matemáticas), cuerpo 11 pt, interlineado holgado. Compilar con Tectonic o XeLaTeX.

\usepackage[spanish,es-noshorthands,es-tabla,es-nodecimaldot]{babel}
\usepackage{libertinus}
\usepackage{microtype}
\usepackage[a4paper,top=2.7cm,bottom=2.7cm,inner=2.8cm,outer=2.4cm,headheight=15pt]{geometry}
\linespread{1.13}
\clubpenalty=10000
\widowpenalty=10000
\displaywidowpenalty=10000
\raggedbottom
\emergencystretch=2.5em
\tolerance=1200
\hbadness=3000

\usepackage{amsmath}
\usepackage{etoolbox}
\usepackage{xcolor}
\usepackage{graphicx}
\usepackage{placeins}
\usepackage{float}
\usepackage{caption}
\usepackage{subcaption}
\usepackage{booktabs,tabularx,array,multirow,makecell}
\usepackage{enumitem}
\usepackage{needspace}
\usepackage{fancyhdr}
\usepackage{titlesec}
\usepackage{listings}
\usepackage{tikz}
\usetikzlibrary{arrows.meta,positioning,shapes.geometric,fit,calc}
\usepackage{algorithm}
\usepackage{algpseudocode}
\usepackage[hidelinks,bookmarksnumbered]{hyperref}

% ── Colores (los mismos de las figuras y de la web) ──
\definecolor{azul}{HTML}{2A78D6}
\definecolor{naranja}{HTML}{EB6834}
\definecolor{aqua}{HTML}{1BAF7A}
\definecolor{amarillo}{HTML}{EDA100}
\definecolor{rosa}{HTML}{E87BA4}
\definecolor{verde}{HTML}{008300}
\definecolor{violeta}{HTML}{6250D6}
\definecolor{rojo}{HTML}{E34948}
\definecolor{gris}{HTML}{52514E}
\definecolor{grisclaro}{HTML}{F1F0EC}
\definecolor{tinta}{HTML}{14203A}

% ── Títulos: cada capítulo empieza en página nueva; los apartados no se quedan huérfanos al pie de la página ──
\titleformat{\chapter}[display]{\normalfont\bfseries\color{tinta}}{\large\scshape\color{gris}Capítulo \thechapter}{0.35em}{\Huge}[\vspace{0.4em}\titlerule]
\titlespacing*{\chapter}{0pt}{-10pt}{26pt}
\titleformat{\section}{\normalfont\Large\bfseries\color{tinta}}{\thesection}{0.7em}{}
\titlespacing*{\section}{0pt}{2.0ex plus 0.4ex}{1.1ex plus 0.2ex}
\titleformat{\subsection}{\normalfont\large\bfseries\color{tinta}}{\thesubsection}{0.7em}{}
\titlespacing*{\subsection}{0pt}{1.7ex plus 0.3ex}{0.8ex plus 0.2ex}
\titleformat{\subsubsection}{\normalfont\normalsize\bfseries\itshape\color{gris}}{}{0pt}{}
\titlespacing*{\subsubsection}{0pt}{1.3ex plus 0.2ex}{0.5ex plus 0.1ex}
% Si queda menos de ~28 % de página, el apartado empieza en la página siguiente (así no se parte por la mitad)
\pretocmd{\chapter}{\FloatBarrier}{}{}
\pretocmd{\section}{\Needspace{0.40\textheight}}{}{}
\pretocmd{\subsection}{\Needspace{0.22\textheight}}{}{}
\pretocmd{\subsubsection}{\Needspace{0.10\textheight}}{}{}

% ── Cabecera y pie ──
\pagestyle{fancy}
\fancyhf{}
\fancyhead[LE]{\small\slshape\leftmark}
\fancyhead[RO]{\small\slshape\rightmark}
\fancyfoot[C]{\small\thepage}
\renewcommand{\headrulewidth}{0.3pt}
\renewcommand{\chaptermark}[1]{\markboth{#1}{}}
\renewcommand{\sectionmark}[1]{\markright{\thesection\ \ #1}}
\fancypagestyle{plain}{\fancyhf{}\fancyfoot[C]{\small\thepage}\renewcommand{\headrulewidth}{0pt}}

% ── Figuras, tablas y listas ──
\captionsetup{font=small,labelfont=bf,labelsep=period,justification=centering,format=plain}
\captionsetup[table]{position=above,skip=6pt}
\captionsetup[figure]{position=below,skip=6pt}
\setlist{itemsep=0.25em,topsep=0.4em,parsep=0pt}
\renewcommand{\arraystretch}{1.18}
\graphicspath{{figuras/}{../web/capturas/}{capturas/}}
\renewcommand{\topfraction}{0.88}\renewcommand{\bottomfraction}{0.7}\renewcommand{\textfraction}{0.12}\renewcommand{\floatpagefraction}{0.65}\setcounter{topnumber}{3}\setcounter{totalnumber}{5}
\setcounter{tocdepth}{2}
\setcounter{secnumdepth}{3}

% ── Pseudocódigo en español ──
\floatname{algorithm}{Algoritmo}
\renewcommand{\listalgorithmname}{Lista de algoritmos}
\algrenewcommand\algorithmicrequire{\textbf{Entrada:}}
\algrenewcommand\algorithmicensure{\textbf{Salida:}}
\algrenewcommand\algorithmicwhile{\textbf{mientras}}
\algrenewcommand\algorithmicdo{\textbf{hacer}}
\algrenewcommand\algorithmicfor{\textbf{para}}
\algrenewcommand\algorithmicforall{\textbf{para cada}}
\algrenewcommand\algorithmicif{\textbf{si}}
\algrenewcommand\algorithmicthen{\textbf{entonces}}
\algrenewcommand\algorithmicelse{\textbf{si no}}
\algrenewcommand\algorithmicrepeat{\textbf{repetir}}
\algrenewcommand\algorithmicuntil{\textbf{hasta que}}
\algrenewcommand\algorithmicreturn{\textbf{devolver}}
\algrenewtext{EndWhile}{\textbf{fin mientras}}
\algrenewtext{EndFor}{\textbf{fin para}}
\algrenewtext{EndIf}{\textbf{fin si}}
\algrenewtext{Else}{\textbf{si no}}
\algrenewtext{ElsIf}[1]{\textbf{si no, si} #1 \textbf{entonces}}
\algrenewtext{Procedure}[2]{\textbf{procedimiento} \textsc{#1}(#2)}
\algrenewtext{EndProcedure}{\textbf{fin procedimiento}}
\algrenewcommand\alglinenumber[1]{\footnotesize\color{gris}#1}
\newcommand{\Comentario}[1]{\hfill\textcolor{gris}{\small\textit{#1}}}

% ── Código y comandos ──
\lstset{basicstyle=\ttfamily\small,breaklines=true,breakatwhitespace=false,columns=fullflexible,keepspaces=true,frame=single,rulecolor=\color{grisclaro},
        backgroundcolor=\color{grisclaro},xleftmargin=0.4em,xrightmargin=0.4em,aboveskip=0.8em,belowskip=0.6em,showstringspaces=false,
        inputencoding=utf8,extendedchars=true,literate={á}{{\'a}}1 {é}{{\'e}}1 {í}{{\'i}}1 {ó}{{\'o}}1 {ú}{{\'u}}1 {ñ}{{\~n}}1 {Á}{{\'A}}1 {É}{{\'E}}1 {Ó}{{\'O}}1 {¿}{{?`}}1 {¡}{{!`}}1 {→}{{$\rightarrow$}}1}
\newcommand{\cod}[1]{\texttt{\small #1}}
\newcommand{\Oh}[1]{\ensuremath{O\!\left(#1\right)}}

% ── Cuadros destacados ──
\newcommand{\recuadro}[2]{\par\vspace{0.6em}\noindent\begin{minipage}{\linewidth}\noindent\colorbox{grisclaro}{\parbox{\dimexpr\linewidth-2\fboxsep}{\small\textbf{#1.} #2}}\end{minipage}\par\vspace{0.6em}}

% ── Datos de portada ──
\newcommand{\titulotrabajo}{El problema del viajante sobre la red de carreteras de España}
\newcommand{\subtitulotrabajo}{Construcción del grafo vial, algoritmos exactos, heurísticos y metaheurísticos, y estudio empírico de su complejidad}
\newcommand{\autorestrabajo}{Raimundo, Pablo, Juan y Jiahao}
\newcommand{\asignatura}{Asignatura MATI \textperiodcentered\ Curso 2026--27}
\newcommand{\fechatrabajo}{Noviembre de 2026}

% Formato de los capítulos de los apéndices (se activa tras ppendix)
\newcommand{\apendicestyle}{\titleformat{\chapter}[display]{\normalfont\bfseries\color{tinta}}{\large\scshape\color{gris}Apéndice \thechapter}{0.35em}{\Huge}[\vspace{0.4em}\titlerule]}
